\documentclass[11pt]{article}

\usepackage{kotex}
\usepackage{amsmath, amssymb, amsthm}
\usepackage{hyperref}
\usepackage{geometry}
\geometry{margin=1in}

\theoremstyle{definition}
\newtheorem{definition}{Definition}[section]
\newtheorem{remark}[definition]{Remark}
\theoremstyle{plain}
\newtheorem{theorem}[definition]{Theorem}
\newtheorem{lemma}[definition]{Lemma}
\newtheorem{corollary}[definition]{Corollary}

\newcommand{\MSO}{\mathrm{MSO}}
\newcommand{\MSOtwo}{\mathrm{MSO}_2}
\newcommand{\lean}[1]{\texttt{#1}}
\newcommand{\leanfile}[1]{\texttt{#1}}

\title{GraphMSO Lean 형식화 해설 II: Nice Tree Decomposition}
\author{}
\date{}

\begin{document}
\maketitle

\tableofcontents

\section{개요}

이 문서는 주어진 tree decomposition을 \emph{nice} tree decomposition으로
normalize하는 부분을 다룬다. 대응하는 참조 문서는
\leanfile{Courcelle/nice\_tree\_decomp.tex}이며, 그 Theorem
\lean{thm:nice-normalization}이 이 문서 전체의 목표다.

\begin{theorem}[목표: \leanfile{nice\_tree\_decomp.tex}, \lean{thm:nice-normalization}]
\label{nice:thm:goal}
$G$의 tree decomposition $(T, \beta)$가 모든 bag에 대해
$|\beta(t)| \le \omega+1$을 만족한다고 하자. $n = |N(T)|$라 쓰면, $G$의
nice tree decomposition $(T^\ast, \beta^\ast)$로서
\[
  |\beta^\ast(x)| \le \omega + 1 \quad (\forall x), \qquad
  |N(T^\ast)| \le (3\omega+5)\,n
\]
을 만족하는 것을 구성하는 algorithm이 존재한다.
\end{theorem}

Lean 형식화는 이 목표를 \emph{두 갈래}로 수행한다. 이 이중 구조가 이 문서를
읽을 때 가장 먼저 잡아야 할 그림이다.

\begin{itemize}
  \item \textbf{Proof-facing 갈래} (\leanfile{normalization.lean}): 입력이
    \S\ref{nice:sec:predicate}의 수학적 \lean{RootedTreeDecomposition}
    (\leanfile{Lean\_Basics.tex}의 \lean{basics:def:treedecomp})이다.
    Root의 선택, child의 열거 등이 모두 \emph{noncomputable}하므로 결과도
    noncomputable하지만, 수학적 진술로서는 가장 일반적이다.
  \item \textbf{Executable 갈래} (\leanfile{execDecomp.lean},
    \leanfile{execNormalization.lean}, \leanfile{execColoring.lean}): 입력이
    rose tree \lean{DecompTree}, 즉 실제 데이터다. 전부 계산 가능하며
    \lean{\#eval}로 돌려 볼 수 있다.
\end{itemize}

두 갈래는 서로를 경유하지 않고 \emph{평행하게} 같은 정리를 증명한다. 둘의
공통 도착점은 \S\ref{nice:sec:inductive}의 constructor-coded tree
\lean{InductiveNiceTree}이다.

\begin{center}
\begin{tabular}{lcl}
  \lean{RootedTreeDecomposition} & $\xrightarrow{\ \lean{normalizeCode}\ }$ & \lean{InductiveNiceTree V ∅} \\
  \lean{DecompTree} (rose tree) & $\xrightarrow{\ \lean{normalizeCode}\ }$ & \lean{InductiveNiceTree V ∅} \\
\end{tabular}
\end{center}

\section{Predicate 방식의 nice decomposition}
\label{nice:sec:predicate}

\emph{파일: \leanfile{GraphMSO/Decomp/nice.lean}}

\emph{참조: \leanfile{Courcelle/nice\_tree\_decomp.tex} Definition 1.1
(Nice tree-decomposition)}

이 파일은 참조 문서의 nice 조건 (N1)--(N4)를 그대로 predicate로 옮긴다.
먼저 node의 모양을 분류하는 predicate들이 정의된다.

\begin{definition}[node 모양 predicate]
\label{nice:def:nodeshapes}
\lean{RootedTreeDecomposition T}와 node $t$에 대해:
\begin{itemize}
  \item \lean{IsLeafNode} $t$ --- $t$는 child가 없다.
  \item \lean{HasAtMostTwoChildren} $t$ --- 어떤 두 node
    \lean{left}, \lean{right}가 존재해서 $t$의 모든 child가 그 둘 중 하나다.
  \item \lean{HasExactlyOneChild} $t$ --- unique child가 존재한다
    (\lean{HasUniqueChild}).
  \item \lean{HasExactlyTwoChildren} $t$ --- 서로 다른 두 child가 있고, 모든
    child가 그 둘 중 하나다.
  \item \lean{IsIntroduceNode} $t$ --- unique child가 있고, 어떤 $v$에 대해
    $B_{\text{child}} = B_t \setminus \{v\}$이며 $v \in B_t$.
  \item \lean{IsForgetNode} $t$ --- unique child가 있고, 어떤 $v$에 대해
    $B_t = B_{\text{child}} \setminus \{v\}$이며 $v \in B_{\text{child}}$.
  \item \lean{IsJoinNode} $t$ --- child가 정확히 둘이고, 두 child bag이
    모두 $B_t$와 같다.
\end{itemize}
\end{definition}

\begin{remark}[\lean{HasAtMostTwoChildren}이 cardinality를 쓰지 않는 이유]
\label{nice:rem:atmosttwo}
``child가 두 개 이하''를 $|\{c : \lean{IsChild}\ t\ c\}| \le 2$로 쓰지 않고
``두 개의 후보 \lean{left}, \lean{right}로 덮인다''로 썼다. 두 후보가 같아도
되므로 child가 $0$개나 $1$개인 경우도 자동으로 포함된다.

이렇게 쓴 이유는 cardinality를 다루려면 child들의 집합에 \lean{Fintype}
instance나 \lean{Set.ncard} 계산이 필요한데, 그 집합이 $t$에 의존하는
subtype이라 다루기 번거롭기 때문이다. 후보 두 개로 덮는 형태는 순수하게
일차 논리적이라 case 분석이 훨씬 쉽다. 실제로 이 선택 덕분에
\lean{hasExactlyTwoChildren\_of\_hasAtMostTwoChildren\_of\_not\_isLeafNode\_of\_not\_hasExactlyOneChild}
(``leaf도 아니고 unary도 아니면 정확히 두 child'')를 순수 논리 조작만으로
증명할 수 있다.
\end{remark}

\begin{definition}[\lean{IsNice}와 \lean{NiceTreeDecomposition}]
\label{nice:def:isnice}
\lean{IsNice T}는 다섯 조건의 conjunction이다: root bag이 $\emptyset$,
모든 leaf bag이 $\emptyset$, 모든 node가 \lean{HasAtMostTwoChildren},
unary node는 introduce이거나 forget, binary node는 join.
\lean{NiceTreeDecomposition}은 \lean{RootedTreeDecomposition}에
\lean{nice : IsNice}를 추가한 structure다.
\end{definition}

\begin{remark}[이 파일은 실제로는 \emph{쓰이지 않는다}]
\label{nice:rem:niceunused}
읽는 사람이 반드시 알아야 할 사실이다. \lean{IsNice}와
\lean{NiceTreeDecomposition}은 이 파일에서 정의된 뒤 프로젝트의 다른 어느
곳에서도 사용되지 않는다. Pipeline 전체가 쓰는 것은
\S\ref{nice:sec:inductive}의 \lean{InductiveNiceTreeDecomposition}이다.

즉 \leanfile{nice.lean}은 \emph{참조 정의(reference definition)}로 봐야 한다.
``nice란 무엇인가''를 참조 문서와 일대일로 대응시켜 기록해 두는 역할이고,
실제 증명은 다음 절의 방식으로 그 조건들을 \emph{정의에 의해} 만족시킨다.
그래서 ``constructor 방식이 predicate 방식의 조건을 만족한다''는 변환 정리는
코드에 존재하지 않는다.

이것은 결함이라기보다 설계 결과다. \S\ref{nice:rem:indexedidea}에서 보듯이
constructor 방식에서는 (N1)--(N4)가 증명할 대상이 아니라 타입 자체가 강제하는
사실이 되기 때문이다. 다만 \leanfile{docs/PLAN.md}에는 두 방식 사이의 변환이
완료된 것처럼 적혀 있으므로, 그 기술은 코드보다 앞서 있는 셈이다.
\end{remark}

\section{Constructor 방식의 nice tree}
\label{nice:sec:inductive}

\emph{파일: \leanfile{GraphMSO/Decomp/nice\_inductive.lean}}

\subsection{핵심 아이디어: bag으로 index된 inductive type}

\begin{definition}[\lean{InductiveNiceTree}]
\label{nice:def:inductivenicetree}
\[
  \lean{InductiveNiceTree}\ V : \lean{Set}\ V \to \lean{Type}
\]
는 root bag으로 \emph{index된} inductive type이며, 네 개의 constructor를 갖는다.
\[
\begin{array}{ll}
  \lean{leaf} & : \lean{InductiveNiceTree}\ V\ \emptyset \\[2pt]
  \lean{introduce}\ v\ (\text{child} : \lean{InductiveNiceTree}\ V\ B)\ (v \notin B)
    & : \lean{InductiveNiceTree}\ V\ (B \cup \{v\}) \\[2pt]
  \lean{forget}\ v\ (\text{child} : \lean{InductiveNiceTree}\ V\ B)\ (v \in B)
    & : \lean{InductiveNiceTree}\ V\ (B \setminus \{v\}) \\[2pt]
  \lean{join}\ (\text{left}, \text{right} : \lean{InductiveNiceTree}\ V\ B)
    & : \lean{InductiveNiceTree}\ V\ B
\end{array}
\]
\end{definition}

\begin{remark}[indexed inductive type이라는 발상]
\label{nice:rem:indexedidea}
이것이 이 파일 전체, 나아가 이 프로젝트의 nice decomposition 취급 방식을
결정하는 아이디어다.

보통이라면 ``tree를 정의하고, 그것이 nice하다는 조건을 predicate로 붙인다''
(\S\ref{nice:sec:predicate}). 여기서는 반대로 \emph{타입의 index에 root bag을
넣어} constructor 자체가 bag의 변화를 지정하게 만들었다. 결과적으로:
\begin{itemize}
  \item \lean{leaf}의 타입이 $\emptyset$으로 index되어 있으므로
    \emph{leaf bag이 비어 있다는 것은 증명할 것이 없다}. (N1)의 절반.
  \item \lean{introduce}는 반드시 정확히 한 vertex를 더하고, \lean{forget}은
    반드시 정확히 한 vertex를 뺀다. 게다가 \lean{fresh : v ∉ bag}과
    \lean{present : v ∈ bag}이 constructor의 인자이므로 ``정확히 하나가
    바뀐다''가 강제된다. (N3).
  \item \lean{join}의 두 child가 \emph{같은} index $B$를 갖도록 타입이 요구하고,
    결과도 $B$다. 즉 세 bag이 모두 같다. (N4).
  \item constructor가 네 개뿐이고 child를 최대 두 개만 받으므로 (N2).
  \item 최종 결과물의 타입을 \lean{InductiveNiceTree V ∅}로 잡으면 root bag이
    비어 있다. (N1)의 나머지 절반.
\end{itemize}
따라서 참조 문서의 Lemma \lean{lem:nice-local} (``구성한 것이 국소적으로
nice하다'')에 해당하는 증명이 Lean에는 \emph{존재하지 않는다}. 타입 검사기가
이미 확인했기 때문이다. \S\ref{nice:rem:niceunused}에서 말한 상황이 이것이다.

대가도 있다. 타입 index가 \emph{구문적으로} 맞아야 하므로,
$B \cup \{v\}$처럼 계산된 표현과 다른 형태로 쓰인 같은 집합을 잇는 명시적
transport가 필요해진다 (\S\ref{nice:def:castroot}).
\end{remark}

\begin{remark}[universe 이야기]
Lean 코드에서 이 타입은 \lean{Type (u+1)}에 산다 ($V : \lean{Type}\ u$일 때).
Index인 \lean{Set V}가 \lean{Type u}이므로 그 위의 inductive family는 한 단계
위 universe에 놓이기 때문이다. 순수하게 기술적인 사항이며 수학적 내용과는
무관하다. 다만 이 때문에 뒤에서 \lean{Node} 타입도 \lean{Type (u+1)}에 살고,
automaton의 state type과 맞추는 지점 (\leanfile{Lean\_Automata.tex})에서
universe를 신경 써야 한다.
\end{remark}

\subsection{Code 안의 위치(position)}

Constructor code는 tree ``전체''를 하나의 값으로 준다. 그 안의 개별 node를
가리키려면 별도의 타입이 필요하다.

\begin{definition}[\lean{Node}, \lean{root}, \lean{nodeBag}, \lean{children}]
\label{nice:def:codenode}
\lean{Node tree}는 tree의 구조에 대한 재귀로 정의된다:
\[
\begin{array}{ll}
  \lean{Node}\ \lean{leaf} &= \lean{PUnit} \\
  \lean{Node}\ (\lean{introduce}\ \_\ c\ \_) &= \lean{Option}\ (\lean{Node}\ c) \\
  \lean{Node}\ (\lean{forget}\ \_\ c\ \_) &= \lean{Option}\ (\lean{Node}\ c) \\
  \lean{Node}\ (\lean{join}\ l\ r) &= \lean{Option}\ (\lean{Node}\ l \oplus \lean{Node}\ r)
\end{array}
\]
즉 \lean{none}이 그 constructor의 root position이고, \lean{some}은 child로
내려간 position이다. \lean{root}, \lean{nodeBag} (각 position의 bag),
\lean{children} (각 position의 child 목록), \lean{depth}가 같은 방식의 재귀로
정의된다.
\end{definition}

\begin{remark}
Position의 개수가 유한하다는 것 (\lean{nodeFintype})도 재귀로 얻는다:
\lean{PUnit}, \lean{Option}, $\oplus$가 모두 유한성을 보존하기 때문이다.
\lean{IsChild parent child}는 ``\lean{child}가 \lean{children tree parent}
목록에 들어 있다''로 정의된다. 즉 \leanfile{Lean\_Basics.tex}의
\lean{basics:def:parent}와 달리 여기서는 child 관계가 \emph{원시 데이터}이고
parent가 파생물이다.
\end{remark}

\subsection{Code가 진짜 tree임을 확인하기}

\begin{definition}[\lean{graph}]
\label{nice:def:codegraph}
\lean{tree.graph}는 \lean{Node tree} 위의 simple graph로, 인접성을
``\lean{IsChild} 관계의 두 orientation 중 하나''로 정의한다.
\end{definition}

\begin{theorem}[\lean{graph\_isTree}]
\label{nice:thm:graphistree}
모든 constructor code에 대해 \lean{tree.graph}는 tree다.
\end{theorem}

\begin{remark}[증명 방식]
Connectivity (\lean{graph\_connected})는 모든 position이 root의 descendant임을
(\lean{root\_reflTransGen\_isChild}) 구조 귀납으로 보이면 된다.

Acyclicity (\lean{graph\_isAcyclic})가 흥미롭다. 가상의 cycle을 잡고 그 위에서
\lean{depth}가 \emph{최대}인 position을 본다. Cycle 위의 점은 서로 다른 이웃을
두 개 가져야 하는데, code graph의 이웃은 parent 아니면 child이고, child는
depth가 더 크다 (\lean{depth\_eq\_add\_one\_of\_isChild}). 최대성 때문에 이웃은
모두 parent여야 하는데, parent는 유일하다 (\lean{parent\_unique}). 모순이다.

핵심 보조 사실 두 개인 \lean{depth\_eq\_add\_one\_of\_isChild}와
\lean{parent\_unique}는 네 constructor $\times$ position의 \lean{none}/\lean{some}
경우를 전부 나열하는 긴 case 분석으로 증명된다. 수학적으로는 자명하지만
Lean에서는 \lean{Option}과 $\oplus$의 constructor를 일일이 벗겨야 해서
분량이 크다.
\end{remark}

또 \lean{graph\_dist\_root\_eq\_depth}: code graph에서 root로부터의 거리가
정의된 \lean{depth}와 일치한다. 이것이 code의 조합 구조와
\leanfile{Lean\_Basics.tex}의 \lean{basics:def:parent} (거리로 정의된 depth)를
연결한다.

\subsection{수학적 decomposition으로의 실현(realization)}

\begin{definition}[\lean{toRootedTreeDecomposition}]
\label{nice:def:torooted}
Constructor code와 \emph{세 개의 bag 공리} (vertex coverage, edge coverage,
occurrence preconnectedness)의 증명을 받아
\lean{RootedTreeDecomposition G}를 만든다. Tree 구조 자체
(\lean{Node}, \lean{graph}, \lean{graph\_isTree}, \lean{root})는 code가 이미
제공하므로 호출자는 bag 공리만 증명하면 된다.
\end{definition}

\begin{theorem}[\lean{toRootedTreeDecomposition\_isChild\_iff}]
\label{nice:thm:childiff}
이렇게 만든 decomposition의 (거리로 정의된) child 관계는 code의
\lean{IsChild}와 정확히 일치한다.
\end{theorem}

\begin{remark}
증명은 양방향이다. ($\Leftarrow$)는 code child가 인접이고 depth를 $1$ 늘리므로
\leanfile{Lean\_Basics.tex}의
\lean{isChild\_of\_adj\_of\_rootDepth\_eq\_add\_one}을 쓴다. ($\Rightarrow$)는
수학적 child가 인접이므로 code child가 두 orientation 중 하나인데, 반대
orientation이면 depth 관계가 모순임을 \lean{omega}로 보인다.
\end{remark}

\begin{definition}[\lean{Realizes}]
\label{nice:def:realizes}
Code \lean{tree}와 수학적 decomposition $T$ 사이의 호환성을 담은 structure다:
position에서 node로 가는 map \lean{realize}, 그것이 bijective라는 증거,
root를 root로 보낸다는 증거, bag을 보존한다는 증거, child 관계를 정확히
대응시킨다는 증거.
\end{definition}

\begin{definition}[\lean{InductiveNiceTreeDecomposition}]
\label{nice:def:indnicedecomp}
\lean{RootedTreeDecomposition G}에 code \lean{tree : InductiveNiceTree V ∅}와
\lean{realization : Realizes tree toRootedTreeDecomposition}을 추가한
structure. 이것이 pipeline 전체가 실제로 사용하는 ``nice decomposition''이다.
\end{definition}

\begin{remark}[왜 \lean{Realizes}를 따로 두는가]
\label{nice:rem:whyrealizes}
\lean{toInductiveNiceTreeDecomposition}으로 만들 때는 \lean{realize}가 그냥
identity이므로 \lean{Realizes}가 자명해 보인다. 그런데도 structure의 field로
둔 이유는, \emph{다른 경로로 만들어진} 수학적 decomposition에 code를 붙이고
싶을 때가 있기 때문이다. Code node와 수학적 node가 서로 다른 타입일 수
있으므로, 둘을 잇는 bijection을 명시적으로 갖고 다녀야 한다.

Root bag이 비어 있다는 사실 (\lean{root\_empty})은 이 구조에서
\emph{따라 나온다}: code의 타입 index가 $\emptyset$이고 \lean{Realizes}가
bag을 보존하기 때문이다. 별도의 field가 아니다.
\end{remark}

\section{Realization을 따라간 성질 전달}
\label{nice:sec:realization}

\emph{파일: \leanfile{GraphMSO/Decomp/realization.lean}}

Code 쪽에도 width와 bag coloring의 개념이 필요하다.

\begin{definition}
\label{nice:def:codewidth}
\lean{tree.HasWidth ω} $:= \forall n.\ |\lean{nodeBag}\ \lean{tree}\ n| \le \omega+1$,
\lean{tree.IsBagColoring color} $:= \forall n.\ \lean{Set.InjOn}\ \lean{color}\ (\lean{nodeBag}\ \lean{tree}\ n)$.
\end{definition}

\begin{theorem}[\lean{hasWidth\_iff}, \lean{isBagColoring\_iff}]
\label{nice:thm:transfer}
\lean{Realizes tree T}가 있으면
\[
  T.\lean{HasWidth}\ \omega \iff \lean{tree}.\lean{HasWidth}\ \omega,
  \qquad
  T.\lean{IsBagColoring}\ c \iff \lean{tree}.\lean{IsBagColoring}\ c.
\]
\end{theorem}

증명은 \lean{realize}가 bijective이고 bag을 보존한다는 것에서 바로 나온다.
$\iff$ 형태로 두었으므로 어느 쪽을 ``진실의 원천''으로 삼아도 된다.

\begin{corollary}[\lean{exists\_bagColoring\_of\_codeHasWidth}]
\label{nice:cor:codecoloring}
Code 쪽 width bound만 주어져도 $\{0,\dots,\omega\}$를 쓰는 code 쪽
bag-injective coloring이 존재한다.
\end{corollary}

증명 경로가 전형적이다: code width $\to$ (전달) $\to$ 수학적 width $\to$
(\leanfile{Lean\_Basics.tex}의 \lean{basics:thm:existsbagcoloring}) $\to$
수학적 coloring $\to$ (전달) $\to$ code coloring. 즉 realization을 두 번
건넌다.

\section{Proof-facing normalization}
\label{nice:sec:normalization}

\emph{파일: \leanfile{GraphMSO/Decomp/normalization.lean} (2487줄)}

\emph{참조: \leanfile{Courcelle/nice\_tree\_decomp.tex} \S Algorithm,
\S Verification}

프로젝트에서 가장 큰 파일이다. 구조는 아래에서 위로 쌓인다.

\subsection{Code에 대한 술어와 크기}

\begin{definition}
\label{nice:def:codepredicates}
\begin{itemize}
  \item \lean{HasBag tree target} --- 어떤 position의 bag이 \lean{target}과 같다.
  \item \lean{VerticesSubset tree allowed} --- 모든 bag이 \lean{allowed}에
    포함된다.
  \item \lean{Occurs tree v} --- $v$가 어떤 bag에 나타난다.
  \item \lean{OccPreconnected tree v} --- $v$를 담는 position들이 code graph에서
    preconnected하다.
  \item \lean{size tree} --- constructor의 개수 (\lean{leaf}는 $1$,
    \lean{introduce}/\lean{forget}은 child$+1$, \lean{join}은 두 child의 합$+1$).
\end{itemize}
\end{definition}

\begin{theorem}[\lean{card\_node\_eq\_size}]
\label{nice:thm:cardsize}
$|\lean{Node}\ \lean{tree}| = \lean{tree}.\lean{size}$.
\end{theorem}

이 등식이 있어서 ``constructor 개수''라는 계산 가능한 양으로 node 개수를 세고,
나중에 정리 \ref{nice:thm:goal}의 $|N(T^\ast)|$ 상계로 옮길 수 있다.

\subsection{\lean{castRoot}: 타입 index 옮기기}

\begin{definition}[\lean{castRoot}]
\label{nice:def:castroot}
$A = B$의 증명과 \lean{InductiveNiceTree V A}를 받아
\lean{InductiveNiceTree V B}를 준다.
\end{definition}

\begin{remark}[왜 필요한가]
\label{nice:rem:castrootwhy}
\S\ref{nice:rem:indexedidea}에서 예고한 대가다. 예를 들어 bag $B$에서
$B \setminus \{v_1\} \setminus \{v_2\} \cup \{w\}$까지 내려간 결과의 타입 index는
그 계산식 그대로다. 이것이 우리가 원하는 표현 $A$와 \emph{집합으로서는} 같아도
Lean이 자동으로 같다고 보지 않으므로, 등식의 증명을 명시적으로 건네
타입을 옮겨야 한다.

그래서 \lean{HasWidth}, \lean{HasBag}, \lean{VerticesSubset}, \lean{Occurs},
\lean{OccPreconnected}, \lean{IsBagColoring} 각각에 대해
``\lean{castRoot}를 거쳐도 변하지 않는다''는 보조정리
(\lean{hasWidth\_castRoot\_iff} 등)가 한 벌 존재한다. 코드에서
\lean{castRoot}가 자주 보이는 것은 수학적 내용이 아니라 순전히 이 사정
때문이므로, 읽을 때는 무시해도 좋다.
\end{remark}

\subsection{한 vertex씩 바꾸는 path}

\emph{참조: \leanfile{nice\_tree\_decomp.tex} Lemma \lean{lem:bag-path}}

\begin{definition}[\lean{forgetList}, \lean{introduceList}]
\label{nice:def:forgetintroducelist}
Vertex 목록을 받아 \lean{forget} (각각 그 시점의 bag에 있어야 함) 또는
\lean{introduce} (각각 그 시점의 bag에 없어야 함)를 차례로 적용한다.
\end{definition}

\begin{definition}[\lean{changeRoot}]
\label{nice:def:changeroot}
유한 bag $A, B$와 root bag이 $B$인 code를 받아, 먼저 $B \setminus A$를 모두
forget하고 다음으로 $A \setminus B$를 모두 introduce하여 root bag이 $A$인
code를 만든다.
\[
  B \;\rightsquigarrow\; A \cap B \;\rightsquigarrow\; A.
\]
\lean{closeToLeaf A} $:= \lean{changeRoot}\ A\ \emptyset\ \lean{leaf}$,
즉 빈 leaf 위에 $A$까지 올라가는 path.
\end{definition}

이것이 참조 문서 Lemma \lean{lem:bag-path}의 형식화다. 그 lemma가 주장하는
네 가지가 각각 정리로 대응된다.

\begin{itemize}
  \item ``모든 bag이 $\omega+1$ 이하'' --- \lean{changeRoot\_hasWidth}.
    삭제 구간의 bag은 $B$의 subset, 삽입 구간의 bag은 $A$의 subset이므로
    양끝 bag의 크기만 가정하면 된다.
  \item ``양끝 bag이 실제로 나타난다'' --- \lean{changeRoot\_hasBag}.
  \item ``edge 개수는 $|A \setminus B| + |B \setminus A|$'' ---
    \lean{changeRoot\_size}. 이것이 뒤의 크기 상계의 원료다.
  \item ``각 vertex의 occurrence는 연결된 부분경로'' ---
    \lean{changeRoot\_occPreconnected}. 참조 문서 증명의 네 가지 경우 분석
    ($v \in A \cap B$, $A \setminus B$, $B \setminus A$, 그 밖)에 대응한다.
\end{itemize}

\begin{remark}[\lean{Finset} 판과 \lean{List} 판이 둘 다 있다]
\label{nice:rem:twopresentations}
같은 내용이 두 번 나온다: \lean{changeRoot}/\lean{closeToLeaf}는
\lean{Finset}을 받고 (\lean{noncomputable}), \lean{changeRootOfList}/
\lean{closeToLeafOfList}는 \lean{List}를 받는다 (계산 가능).
보조정리 표면도 동일하게 갖춰져 있다.

이유는 \S\ref{nice:sec:normalization}의 proof-facing 갈래는 수학적 bag
(\lean{Set V})에서 출발하므로 \lean{Set.toFinset}을 거쳐야 하는데 이것이
noncomputable이고, \S\ref{nice:sec:exec}의 executable 갈래는 처음부터 raw
list를 갖고 있어 그럴 필요가 없기 때문이다. \lean{List} 판은 중복 원소를
\lean{List.dedup}으로 처리한다.
\end{remark}

\subsection{여러 branch 합치기}

\begin{definition}[\lean{joinNonempty}]
\label{nice:def:joinnonempty}
같은 bag을 갖는 code 하나와 목록을 받아 오른쪽 결합으로 \lean{join}을 반복
적용한다. 참조 문서 Step 5의 ``right-associated binary join tree''.
\end{definition}

크기는 \lean{joinNonempty\_size}: 분기들의 크기 합에 분기 개수 $-1$개의 join
node가 더해진다. Occurrence는 \lean{joinNonempty\_occurs\_iff} (어떤 분기에
나타나면 된다)와 \lean{joinNonempty\_occPreconnected}이 다룬다.

\begin{remark}[\lean{join}의 occurrence 연결성이 까다로운 이유]
\label{nice:rem:joinconn}
\lean{join left right}에서 $v$의 occurrence가 연결되려면, $v$가 양쪽 subtree에
모두 나타나는 경우 두 쪽을 잇는 경로가 root를 지나야 하고, 그러려면
\emph{root bag에 $v$가 있어야} 한다. 이것은 저절로 성립하지 않고 가정으로
들어간다 (\lean{join\_occPreconnected}의 hypothesis). 실제 사용처에서 이
가정을 공급하는 것이 tree decomposition의 running intersection 공리이며,
그 연결 고리가 아래 \lean{HasVertexBelow} 관련 보조정리다.
\end{remark}

\subsection{입력 decomposition에 대한 재귀}

\begin{definition}[\lean{HasVertexBelow}와 그 보조정리]
\label{nice:def:hasvertexbelow}
\lean{HasVertexBelow T t v}는 $t$의 subtree 어딘가의 bag에 $v$가 있다는 뜻이다.
두 개의 핵심 보조정리:
\begin{itemize}
  \item \lean{mem\_child\_bag\_of\_mem\_parent\_of\_hasVertexBelow}: $v$가
    parent bag에 있고 child subtree 아래에도 나타나면, $v$는 child bag에도 있다.
  \item \lean{mem\_parent\_bag\_of\_hasVertexBelow\_two\_children}: $v$가 서로
    다른 두 child subtree 아래에 나타나면 $v$는 parent bag에 있다.
\end{itemize}
\end{definition}

이 둘이 running intersection 공리를 ``재귀에서 쓸 수 있는 국소적 형태''로
바꾼 것이며, \S\ref{nice:rem:joinconn}의 가정을 공급한다.

\begin{definition}[well-founded 재귀 준비]
\label{nice:def:wellfounded}
\lean{childList T t} (child들의 목록), \lean{subtreeFinset T t} (subtree의
node 집합), \lean{subtreeCard T t} (그 크기), 그리고 child 관계가
well-founded임을 보이는 \lean{childWellFounded}.
\end{definition}

\begin{remark}[왜 well-founded 재귀인가]
\label{nice:rem:wfrecursion}
\lean{DecompTree} (\S\ref{nice:sec:exec})와 달리 \lean{RootedTreeDecomposition}은
tree ``데이터''가 아니라 그래프와 공리의 묶음이다. Child 관계는 root로부터의
거리로 \emph{파생된} 관계이므로 Lean이 구조적 재귀를 자동으로 받아들이지
않는다. 그래서 ``child로 내려가면 depth가 커지는데 depth는 node 개수보다 작다''
(\lean{rootDepth\_lt\_card}, \lean{childMeasure\_lt})를 이용해 well-founded
관계를 만들고, \lean{WellFounded.fix}로 재귀를 정의한다.

이것이 proof-facing 갈래가 executable 갈래보다 훨씬 무거운 근본 이유다.
\end{remark}

또 \lean{subtreeCard\_eq\_one\_add\_sum\_childList}: subtree의 node 개수는
$1$ 더하기 child subtree들의 개수 합이다. Child subtree들이 서로소라는
사실 (\lean{subtreeFinset\_children\_pairwiseDisjoint})이 필요하다.

\subsection{Normalization 본체}

\begin{definition}[\lean{normalizeCodeAt}]
\label{nice:def:normalizecodeat}
입력 decomposition $T$와 node $t$에 대해, root bag이 $B_t$인 code를 만든다:
\begin{itemize}
  \item $t$에 child가 없으면 \lean{closeToLeaf} $B_t$, 즉 $B_t$에서 빈 leaf까지
    내려가는 path. (참조 문서 Step 3.)
  \item child들이 있으면, 각 child $s$에 대해 재귀 결과 위에
    \lean{changeRoot} $B_t\ B_s$를 씌워 root bag을 $B_t$로 맞춘 branch를 만들고,
    그 branch들을 \lean{joinNonempty}로 합친다. (Step 4, 5.)
\end{itemize}
\end{definition}

\begin{definition}[\lean{normalizeCode}, \lean{normalize}]
\label{nice:def:normalizecode}
\[
  \lean{normalizeCode}\ T := \lean{changeRoot}\ \emptyset\ B_{\text{root}}\
    (\lean{normalizeCodeAt}\ T\ \lean{root}),
\]
즉 root 위에 빈 bag까지 올라가는 path를 얹는다 (Step 6). 결과의 타입은
\lean{InductiveNiceTree V ∅}.

\lean{normalize T}는 여기에 세 개의 bag 공리를 증명해 붙여
\lean{InductiveNiceTreeDecomposition}을 만든 것이다.
\end{definition}

\subsection{정확성}

참조 문서의 Verification 절과 다음과 같이 대응한다.

\begin{itemize}
  \item Lemma \lean{lem:nice-local} --- 대응하는 Lean 증명 \emph{없음}.
    \S\ref{nice:rem:indexedidea} 참조.
  \item Lemma \lean{lem:original-bags-survive} --- \lean{normalizeCode\_hasBag}:
    원래 decomposition의 모든 bag이 결과에 나타난다.
  \item Lemma \lean{lem:decomp-axioms} --- \lean{normalize}의 세 인자.
    Vertex/edge coverage는 \lean{normalizeCode\_hasBag}에서 바로 나오고,
    running intersection은 \lean{normalizeCode\_occPreconnected}가 준다.
  \item Width 보존 --- \lean{normalizeCode\_hasWidth},
    \lean{normalize\_hasWidth}.
  \item 크기 상계 --- \lean{normalizeCode\_size\_le},
    \lean{normalize\_node\_bound}.
\end{itemize}

\begin{theorem}[\lean{normalize\_node\_bound}]
\label{nice:thm:nodebound}
$T$가 width $\omega$ 이하이면
\[
  |\lean{Node}(\lean{normalize}\ T)| \;\le\; (3\omega+5)\cdot|\lean{Node}\ T|.
\]
\end{theorem}

\begin{remark}[크기 상계의 증명: amortization]
\label{nice:rem:amortization}
이 증명의 설계가 이 파일에서 가장 눈여겨볼 부분이다. 자연스러운 귀납 가설
``$\lean{size}(\lean{normalizeCodeAt}\ t) \le (3\omega+5)\cdot\lean{subtreeCard}\ t$''
는 \emph{통하지 않는다}. $t$를 parent에 붙일 때 \lean{changeRoot}가 최대
$2(\omega+1)$개의 node를 더 쓰는데, 그 비용을 어디서도 감당하지 못하기
때문이다.

그래서 Lean 증명은 상수를 미리 떼어 놓은 강화된 명제를 귀납 가설로 쓴다
(\lean{normalizeCodeAt\_size\_add\_le}):
\[
  \lean{size}(\lean{normalizeCodeAt}\ t) + (2\omega+3)
    \;\le\; (3\omega+5)\cdot \lean{subtreeCard}\ t.
\]
좌변의 $+(2\omega+3)$이 ``$t$에서 부모로 올라가는 edge가 쓸 예산''을 미리
지불해 둔 것이다. 이렇게 하면 귀납 단계에서 각 child의 잉여분을 그 child를
붙이는 \lean{changeRoot} 비용에 정확히 충당할 수 있고, join node 하나
($+1$)와 자기 자신의 몫이 남는다.

이것은 참조 문서 증명의 ``각 original edge가 최대 $2(\omega+1)$개의 새 node를
기여한다''는 전역 세기를 귀납 가능한 국소 형태로 바꾼 것이다. 전역 세기를
그대로 형식화하려면 결과 tree의 node를 원래 edge/leaf/join으로 분류하는
전단사를 만들어야 하는데, amortization은 그 작업을 통째로 피한다.

마지막으로 \lean{normalizeCode\_size\_le}가 root 위의 path
($\le \omega+1$개)를 더하고 $\lean{subtreeCard}(\lean{root}) = |N(T)|$를 써서
정리 \ref{nice:thm:nodebound}를 얻는다.
\end{remark}

무향 입력을 위해 \lean{TreeDecomposition.rooted} (아무 node나 root로 선택)와
\lean{TreeDecomposition.normalize}가 마지막에 붙는다.

\section{Executable 갈래}
\label{nice:sec:exec}

\subsection{Rose tree 표현}

\emph{파일: \leanfile{GraphMSO/Decomp/execDecomp.lean}}

\begin{definition}[\lean{DecompTree}]
\label{nice:def:decomptree}
\[
  \lean{DecompTree}\ V ::= \lean{node}\ (\lean{bag} : \lean{List}\ V)\ (\lean{children} : \lean{List}\ (\lean{DecompTree}\ V)).
\]
Bag은 raw list이며 중복이 있어도 무방하다 (모든 소비자가
\lean{List.toFinset}으로 재는다).
\end{definition}

\begin{remark}[\lean{TreeDecomposition}과의 차이]
\label{nice:rem:rosevsmath}
차이가 본질적이다.
\begin{itemize}
  \item Root, tree 모양, child 목록이 전부 \emph{데이터에 내장}되어 있다.
    따라서 tree 공리도, root 선택도, child 열거도 필요 없고, algorithm이
    \emph{구조적 재귀}로 내려갈 수 있다 (\S\ref{nice:rem:wfrecursion}과 대조).
  \item 그래프 $G$에 대한 유효성은 데이터에서 분리되어 별도 predicate
    \lean{IsDecompFor}로 주어진다. 즉 \emph{계산 함수는 raw 데이터를 먹고,
    정확성 정리만 유효성을 먹는다}. 이 분리가 executable 갈래 설계의 핵심
    원칙이며 \leanfile{Lean\_Algorithm.tex}에서도 반복된다.
\end{itemize}
\end{remark}

\begin{definition}[\lean{RunningIntersection}과 \lean{IsDecompFor}]
\label{nice:def:isdecompfor}
\lean{RunningIntersection}은 inductive predicate로, 각 node에서
\begin{itemize}
  \item \lean{hdown}: node bag의 vertex가 어떤 child subtree에 나타나면 그
    child의 root bag에도 있다;
  \item \lean{hpair}: 서로 다른 두 child subtree에 함께 나타나는 vertex는 node
    bag에 있다;
  \item \lean{hchildren}: 모든 child가 재귀적으로 같은 성질을 갖는다.
\end{itemize}
\lean{IsDecompFor t G}는 vertex coverage, edge coverage,
\lean{RunningIntersection}의 세 field를 갖는 structure다.
\end{definition}

\begin{remark}[connectivity 공리의 국소화]
\leanfile{Lean\_Basics.tex}의 \lean{basics:def:treedecomp}는 running
intersection을 ``$\mathrm{BAGS}(v)$가 preconnected''라는 \emph{전역} 조건으로
쓴다. 여기서는 그것을 rooted tree에서의 두 국소 조건으로 바꿨다. 이 두
조건이 \S\ref{nice:def:hasvertexbelow}의 두 보조정리와 정확히 같은 내용이며,
바로 그렇기 때문에 구조적 재귀 도중에 쓸 수 있다. 검증할 때도 국소 조건이
훨씬 편하다.
\end{remark}

\subsection{계산 가능한 normalization}

\emph{파일: \leanfile{GraphMSO/Decomp/execNormalization.lean}}

\begin{definition}[\lean{normalizeAux}, \lean{normalizeCode}]
\label{nice:def:execnormalize}
\[
  \lean{normalizeAux} : (t : \lean{DecompTree}\ V) \to \lean{InductiveNiceTree}\ V\ (t.\lean{rootBag}.\lean{toFinset}),
\]
child가 없으면 \lean{closeToLeafOfList}, 있으면 각 child를
\lean{changeRootOfList}로 부모 bag에 맞춘 뒤 \lean{joinNonempty}.
\lean{normalizeCode t}는 그 위에 빈 root path를 얹은 것이다.
\end{definition}

\S\ref{nice:def:normalizecodeat}와 문장이 거의 같지만 정의 방식이 다르다:
여기서는 \lean{decreasing\_by}로 child의 크기가 줄어듦만 보이면 되는
\emph{구조적} 재귀다. \lean{WellFounded.fix}도, \lean{subtreeCard}도 필요 없다.

정확성 정리들이 proof-facing 판과 일대일로 대응한다:
\lean{normalizeCode\_hasWidth}, \lean{normalizeCode\_hasBag},
\lean{normalizeCode\_occurs\_iff}, \lean{normalizeCode\_occPreconnected},
\lean{normalizeCode\_isBagColoring}, 그리고 크기 상계
\lean{normalizeCode\_size\_le}:
\[
  \lean{size}(\lean{normalizeCode}\ t) \le (3\omega+5)\cdot \lean{size}\ t.
\]
증명 전략(amortization, \S\ref{nice:rem:amortization})도 동일하다
(\lean{normalizeAux\_size\_add\_le}).

\begin{definition}[\lean{DecompTree.normalize}]
\label{nice:def:execnormalizecert}
유효성 증명 \lean{h : t.IsDecompFor G}를 받아
\lean{InductiveNiceTreeDecomposition}을 만든다.
\end{definition}

\begin{remark}[무엇이 계산 가능하고 무엇이 아닌가]
\label{nice:rem:whatcomputes}
\lean{normalize}는 \lean{noncomputable}로 선언되어 있다. 그러나 실제로
계산 불가능한 것은 \emph{증명서 층}뿐이다: 그 안의 code
\lean{t.normalizeCode}는 완전히 계산 가능하고, \lean{normalize\_tree}가
\lean{(t.normalize h).tree = t.normalizeCode}를 \lean{rfl}로 확인해 준다.

실행 가능한 pipeline (\leanfile{Lean\_Algorithm.tex})은 그래서
\lean{normalize}가 아니라 \lean{normalizeCode}를 직접 소비하고,
\lean{normalize}는 정리를 서술할 때만 등장한다. 이 ``계산 가능한 핵심 +
noncomputable 증명서'' 패턴이 프로젝트 전체에서 반복된다.
\end{remark}

\subsection{계산 가능한 greedy bag coloring}

\emph{파일: \leanfile{GraphMSO/Decomp/execColoring.lean}}

\leanfile{Lean\_Basics.tex}의 \lean{basics:thm:existsbagcoloring}은 coloring의
\emph{존재}만 주므로 실행할 수 없다. 이 파일이 그것을 구체적인 함수로 대체한다.

\begin{definition}[\lean{greedyColoring}]
\label{nice:def:greedycoloring}
Rose tree를 root부터 훑으면서, 각 vertex를 \emph{처음 만나는 bag} (즉 그
vertex의 최상위 bag)에서 그 bag의 이미 색칠된 원소들이 쓰지 않은 첫 번째 색으로
칠한다.

구성 요소: \lean{freshColor} ($\lean{Fin}(\omega+1)$에서 \lean{used}에 없는 첫
색), \lean{usedColors}, \lean{assignVertex} (한 vertex 처리),
\lean{colorBag} (한 bag 처리), \lean{greedyAux} (누산기를 들고 tree를 훑는
재귀).
\end{definition}

\begin{remark}[누산기의 정체]
\label{nice:rem:accumulator}
\lean{greedyAux}가 들고 다니는 누산기는 쌍
\[
  (\lean{f} : V \to \lean{Fin}(\omega+1),\ \ \lean{assigned} : \lean{List}\ V)
\]
이다. 즉 ``현재까지의 색칠 함수''와 ``이미 색이 확정된 vertex 목록''이다.
색칠 함수는 처음부터 전역 함수이고 (미할당 vertex에는 의미 없는 값이 들어
있다), \lean{assigned} 목록이 어느 값이 진짜인지를 기록한다. Lean에서 부분
함수를 다루는 흔한 방식이다.
\end{remark}

\begin{theorem}[\lean{greedyColoring\_isBagColoring}]
\label{nice:thm:greedycorrect}
$t$가 width $\omega$ 이하이고 \lean{RunningIntersection}을 만족하면
\lean{greedyColoring t ω}는 $t$의 모든 bag 위에서 injective하다.
\end{theorem}

\begin{remark}[증명 방식: 명시적 invariant]
\label{nice:rem:greedyinvariant}
증명은 \lean{greedyAux}의 한 번 실행에 대한 사전/사후 조건을 명시적으로
적어 두고 (\lean{GreedyPre}, \lean{GreedyPost}) 그것을 구조적 귀납으로 옮기는
방식이다.

\emph{사전 조건} \lean{GreedyPre}: subtree에 나타나면서 이미 색이 확정된
vertex는 그 subtree의 root bag 안에 있고, 그런 vertex들 위에서 현재 색칠이
injective하다.

\emph{사후 조건} \lean{GreedyPost}: 기존에 확정된 값은 바뀌지 않는다;
새로 확정되는 것은 정확히 그 subtree에 나타나는 vertex들이다; 결과 색칠은
그 subtree의 \emph{모든} bag 위에서 injective하다.

핵심은 사전 조건이 유지된다는 부분이며, 여기서 running intersection이 쓰인다:
현재 bag에서 이미 색칠된 원소는 조상 bag에서 칠해졌거나 앞선 형제 subtree에서
칠해진 것인데, 두 경우 모두 \lean{hdown}/\lean{hpair}에 의해 현재 bag의
부모 bag에 들어 있다. 그 bag의 크기가 $\omega+1$ 이하이므로 막히는 색은 최대
$\omega$개이고, \lean{freshColor\_notMem}이 남는 색의 존재를 보장한다.

이 논증이 \leanfile{Lean\_Basics.tex}의 \lean{basics:rem:greedyproof}에서
쓴 ``최상위 bag'' 논증의 계산 가능한 판본이다. 그쪽은 유한집합에 대한 강귀납법과
$\mathrm{BAGS}$ 볼록성을 썼고, 이쪽은 tree를 실제로 훑는 순서를 이용해 같은
사실을 얻는다.
\end{remark}

\begin{theorem}[\lean{normalizeCode\_greedyColoring\_isBagColoring}]
\label{nice:thm:greedytransfer}
같은 가정 아래, \lean{greedyColoring t ω}는 \emph{normalize한 code}
\lean{t.normalizeCode}의 bag coloring이기도 하다.
\end{theorem}

이것이 실제로 쓰이는 형태다. \lean{normalizeCode\_isBagColoring}
(\S\ref{nice:def:execnormalize}) 을 통해 전달되며, 그 정리는
\lean{changeRootOfList}, \lean{closeToLeafOfList}, \lean{joinNonempty} 각
구성 요소가 \lean{IsBagColoring}을 보존한다는 보조정리들 위에 세워져 있다.
새로 생기는 bag은 모두 원래 bag의 subset이므로 injectivity가 유지된다는 것이
그 내용이다.

\begin{remark}[색의 개수]
이 coloring은 $\omega+1$개의 색을 쓴다. 즉 decomposition의 width에 맞춘
최적 개수다. 대안으로 \leanfile{Lean\_Algorithm.tex}의 \lean{checkFin}은
vertex가 $\lean{Fin}\ n$일 때 identity를 색칠로 쓰는데, 그 경우 색이 $n+1$개
필요하다. Alphabet 크기가 색 개수에 지수적으로 의존하므로 이 차이는 실제
성능에서 결정적이다.
\end{remark}

\section{요약: 정리 \ref{nice:thm:goal}와의 대응}
\label{nice:sec:summary}

\begin{center}
\begin{tabular}{p{0.30\textwidth}p{0.30\textwidth}p{0.30\textwidth}}
\hline
\leanfile{nice\_tree\_decomp.tex} & proof-facing & executable \\
\hline
Definition 1.1 (N1)--(N4) &
  \multicolumn{2}{p{0.62\textwidth}}{\lean{InductiveNiceTree}의 타입이 강제
  (\S\ref{nice:rem:indexedidea}); \lean{IsNice}는 참조용
  (\S\ref{nice:rem:niceunused})} \\
Lemma \lean{lem:bag-path} &
  \lean{changeRoot} 계열 & \lean{changeRootOfList} 계열 \\
Algorithm Step 1--6 &
  \lean{normalizeCodeAt}, \lean{normalizeCode} &
  \lean{normalizeAux}, \lean{normalizeCode} \\
Lemma \lean{lem:nice-local} &
  \multicolumn{2}{p{0.62\textwidth}}{증명 불필요} \\
Lemma \lean{lem:original-bags-survive} &
  \lean{normalizeCode\_hasBag} & \lean{normalizeCode\_hasBag} \\
Lemma \lean{lem:decomp-axioms} &
  \lean{normalize} & \lean{normalize} \\
Width 보존 &
  \lean{normalize\_hasWidth} & \lean{normalize\_hasWidth} \\
$|N(T^\ast)| \le (3\omega+5)n$ &
  \lean{normalize\_node\_bound} & \lean{normalizeCode\_size\_le} \\
Bag coloring &
  \lean{basics:thm:existsbagcoloring} (존재만) &
  \lean{greedyColoring} (구체적) \\
\hline
\end{tabular}
\end{center}

\begin{remark}[형식화되지 않은 부분: 실행 시간]
정리 \ref{nice:thm:goal}의 ``$O(\omega n)$ 시간'' 주장은 이 문서 범위에서
형식화되지 않는다. Node 개수 상계 $(3\omega+5)n$만 증명된다. 시간에 대한
주장은 \leanfile{Lean\_Basics.tex}의 \lean{basics:def:costed} 비용 계층을 통해
\leanfile{Lean\_Algorithm.tex}에서 \emph{추상적 연산 횟수}로만 다루어지며,
거기서도 실제 실행 시간에 대한 주장은 아니다
(\lean{basics:rem:costmeaning}).
\end{remark}

\end{document}
